\SourceLeaf{012}Quippe quaestio haec per rei naturam aliquid vagi implicat, quod limitibus circumscribi nisi per principium aliquatenus arbitrarium nequit. Determinatio alicuius quantitatis per observationem errori maiori minorive obnoxiam, haud inepte comparatur ludo, in quo solae iacturae, lucra nulla, dum quilibet error metuendus iacturae affinis est. Talis ludi dispendium aestimatur e iactura probabili, puta ex aggregato productorum singularum iacturarum possibilium in probabilitates respectivas. Quantae vero iacturae quemlibet observationis errorem aequiparare conveniat, neutiquam per se clarum est; quin potius haec determinatio aliqua ex parte ab arbitrio nostro pendet. Iacturam ipsi errori aequalem statuere manifesto non licet; si enim errores positivi pro iacturis acciperentur, negativi lucra repraesentare deberent. Magnitudo iacturae potius per talem erroris functionem exprimi debet, quae natura sua semper fit positiva. Qualium functionum quum varietas sit infinita, simplicissima, quae hac proprietate gaudet, prae ceteris eligenda videtur, quae absque lite est quadratum: hoc pacto principium supra prolatum prodit.

Ill. \textsc{Laplace} simili quidem modo rem consideravit, sed errorem ipsum semper positive acceptum tamquam iacturae mensuram adoptavit. At ni fallimur haecce ratio saltem non minus arbitraria est quam nostra: utrum enim error duplex aeque tolerabilis putetur quam simplex bis repetitus, an aegrius, et proin utrum magis conveniat, errori duplici momentum duplex tantum, an maius, tribuere, quaestio est neque per se clara, neque demonstrationibus mathematicis decidenda, sed libero tantum arbitrio remittenda. Praeterea negari non potest, ista ratione continuitatem laedi: et propter hanc ipsam caussam modus ille tractationi analyticae magis refragatur, dum ea, ad quae principium nostrum perducit, mira tum simplicitate tum generalitate commendantur.

\Article{7}
Statuendo valorem integralis \(\int xx\varphi x.\mathrm{d}x\) ab \(x=-\infty\) usque ad \(x=+\infty\) extensi \(=mm\), quantitatem \(m\) vocabimus \textit{errorem medium metuendum}, sive simpliciter \textit{errorem medium} observationum, quarum errores indefiniti \(x\) habent probabilitatem relativam \(\varphi x\). Denominationem illam non ad observationes immediatas limitabimus, sed etiam ad determinationes qualescunque ex observationibus derivatas extendemus. Probe autem cavendum est, ne error medius confundatur cum medio arithmetico omnium errorum, de quo in art. 5 locuti sumus.

\SourceLeaf{013}Ubi plura observationum genera, seu plures determinationes ex observationibus petitae, quibus haud eadem praecisio concedenda est, comparantur. \textit{pondus} earum relativum nobis erit quantitas ipsi \(mm\) reciproce proportionalis, dum \textit{praecisio} simpliciter ipsi \(m\) reciproce proportionalis habetur. Quo igitur pondus per numerum exprimi possit, pondus certi observationum generis pro unitate acceptum esse debet.

\Article{8}
Si observationum errores partem constantem implicant, hanc auferendo error medius minuitur, pondus et praecisio augentur. Retinendo signa art. 5, designandoque per \(m'\) errorem medium observationum correctarum, erit
\[
\begin{aligned}
m'm'&=\int x'x'\varphi'x'.\mathrm{d}x'=\int(x-k)^2\varphi x.\mathrm{d}x\\
&=\int xx\varphi x.\mathrm{d}x-2k\int x\varphi x.\mathrm{d}x+kk\int\varphi x.\mathrm{d}x\\
&=mm-2kk+kk=mm-kk.
\end{aligned}
\]
Si autem loco partis constantis veri \(k\) quantitas alia \(l\) ab observationibus ablata esset, quadratum erroris medii novi evaderet \(=mm-2kl+ll=m'm'+(l-k)^2\).

\Article{9}
Denotante \(\lambda\) coëfficientem determinatum, atque \(\mu\) valorem integralis \(\int\varphi x.\mathrm{d}x\) ab \(x=-\lambda m\) usque ad \(x=+\lambda m\), erit \(\mu\) probabilitas, quod error alicuius observationis sit minor quam \(\lambda m\) (sine respectu signi), nec non \(1-\mu\) probabilitas erroris maioris quam \(\lambda m\). Si itaque valor \(\mu=\tfrac12\) respondet valori \(\lambda m=\rho\), error aeque facile infra \(\rho\) quam supra \(\rho\) cadere potest, quocirca \(\rho\) commode dici potest error \textit{probabilis}. Relatio quantitatum \(\lambda\), \(\mu\) manifesto pendet ab indole functionis \(\varphi x\), quae plerumque incognita est. Operae itaque pretium erit, istam relationem pro quibusdam casibus specialibus propius considerare.

I. Si limites omnium errorum possibilium sunt \(-a\) et \(+a\), omnesque errores intra hos limites aeque probabiles, erit \(\varphi x\) inter limites \(x=-a\) et \(x=+a\) constans, et proin \(=\tfrac{1}{2a}\). Hinc \(m=a\BareRoot{\tfrac13}\), nec non \(\mu=\lambda\BareRoot{\tfrac13}\), quamdiu \(\lambda\) non maior quam \(\BareRoot{3}\); denique \(\rho=m\BareRoot{\tfrac34}=0{,}8660254m\), probabilitasque, quod error prodeat errore medio non maior, erit \(=\BareRoot{\tfrac13}=0{,}5773503\).

II. Si ut antea \(-a\) et \(+a\) sunt errorum possibilium limites, errorumque ipsorum probabilitas inde ab errore \(0\) utrimque in progressione arithmetica decrescere supponitur, erit
\SourceLeaf{014}
\[
\begin{aligned}
\varphi x&=\frac{a-x}{aa},\quad\text{pro valoribus ipsius }x\text{ inter }0\text{ et }+a\\[3pt]
\varphi x&=\frac{a+x}{aa},\quad\text{pro valoribus ipsius }x\text{ inter }0\text{ et }-a
\end{aligned}
\]
Hinc deducitur \(m=a\BareRoot{\tfrac16}\), \(\mu=\lambda\BareRoot{\tfrac23}-\tfrac16\lambda\lambda\), quamdiu \(\lambda\) est inter \(0\) et \(\BareRoot{6}\); denique \(\lambda=\BareRoot{6}-\BareRoot{(6-6\mu)}\), quamdiu \(\mu\) inter \(0\) et \(1\), et proin
\[
\rho=m(\BareRoot{6}-\BareRoot{3})=0{,}7174389m
\]
Probabilitas erroris medium non superantis erit in hoc casu
\[
=\BareRoot{\tfrac32}-\tfrac16=0{,}6498299
\]
III. Si functionem \(\varphi x\) proportionalem statuimus huic \(e^{-\frac{xx}{hh}}\) (quod quidem in rerum natura proxime tantum verum esse potest), esse debebit
\[
\varphi x=\frac{e^{-\frac{xx}{hh}}}{h\BareRoot{\pi}}
\]
denotante \(\pi\) semiperipheriam circuli pro radio \(1\), unde porro deducimus
\[
m=h\BareRoot{\tfrac12}
\]
(V. \textit{Disquis. generales circa seriem infinitam} etc. art. 28). Porro si valor integralis
\[
\frac{2}{\BareRoot{\pi}}\int e^{-zz}\mathrm{d}z
\]
a \(z=0\) inchoati denotatur per \(\Theta z\), erit
\[
\mu=\Theta(\lambda\BareRoot{\tfrac12})
\]
\Needspace{14\baselineskip}
Tabula sequens exhibet aliquot valores huius quantitatis:
\begin{center}
\begin{tabular}{r|l}
\multicolumn{1}{c|}{\(\lambda\)}&\multicolumn{1}{c}{\(\mu\)}\\\hline
\(0{,}6744897\)&\(0{,}5\)\\
\(0{,}8416213\)&\(0{,}6\)\\
\(1{,}0000000\)&\(0{,}68268\mathord{\cdot}95\)\\
\(1{,}0364334\)&\(0{,}7\)\\
\(1{,}2815517\)&\(0{,}8\)\\
\(1{,}6448537\)&\(0{,}9\)\\
\(2{,}5758293\)&\(0{,}99\)\\
\(3{,}2918301\)&\(0{,}999\)\\
\(3{,}8905940\)&\(0{,}9999\)\\
\(\infty\)&\(1\)
\end{tabular}
\end{center}

\SourceLeaf{015}\Article{10}
Quamquam relatio inter \(\lambda\) et \(\mu\) ab indole functionis \(\varphi x\) pendet, tamen quaedam generalia stabilire licet. Scilicet qualiscunque sit haec functio, si modo ita est comparata, ut ipsius valor, crescente valore absoluto ipsius \(x\), semper decrescat, vel saltem non crescat, certo erit
\begin{quote}
\(\lambda\) minor vel saltem non maior quam \(\mu\BareRoot{3}\), quoties \(\mu\) est minor quam \(\tfrac23\);

\(\lambda\) non maior quam \(\frac{2}{3\BareRoot{(1-\mu)}}\), quoties \(\mu\) est maior quam \(\tfrac23\).
\end{quote}
Pro \(\mu=\tfrac23\) uterque limes coincidit, puta \(\lambda\) nequit esse maior quam \(\BareRoot{\tfrac43}\).

Ut hoc insigne theorema demonstremus, denotemus per \(y\) valorem integralis \(\int\varphi z.\mathrm{d}z\) a \(z=-x\) usque ad \(z=+x\) extensi, quo pacto \(y\) erit probabilitas, quod error aliquis contentus sit intra limites \(-x\) et \(+x\). Porro statuamus
\[
x=\psi y,\quad \mathrm{d}\psi y=\psi'y.\mathrm{d}y,\quad \mathrm{d}\psi'y=\psi''y.\mathrm{d}y
\]
Erit itaque \(\psi0=0\), nec non
\[
\psi'y=\frac{1}{\varphi x+\varphi(-x)}
\]
quare per hyp. \(\psi'y\) ab \(y=0\) usque ad \(y=1\) semper crescet, saltem nullibi decrescet, sive, quod idem est, valor ipsius \(\psi''y\) semper erit positivus, vel saltem non negativus. Porro habemus \(\mathrm{d}.y\psi'y=\psi'y.\mathrm{d}y+y\psi''y.\mathrm{d}y\), adeoque
\[
y\psi'y-\psi y=\int y\psi''y.\mathrm{d}y
\]
integratione ab \(y=0\) inchoata. Valor expressionis \(y\psi'y-\psi y\) itaque semper erit quantitas positiva, saltem non negativa, adeoque
\[
1-\frac{\psi y}{y\psi'y}
\]
quantitas positiva unitate minor. Sit \(f\) eius valor pro \(y=\mu\), i. e. quum habeatur \(\psi\mu=\lambda m\), sit
\[
f=1-\frac{\lambda m}{\mu\psi'\mu}\quad\text{sive}\quad\psi'\mu=\frac{\lambda m}{(1-f)\mu}
\]
His ita praeparatis, consideremus functionem ipsius \(y\) hanc
\[
\frac{\lambda m}{(1-f)\mu}(y-\mu f)
\]
quam statuemus \(=Fy\), nec non \(\mathrm{d}Fy=F'y.\mathrm{d}y\). Perspicuum est, fieri
